\documentclass[10pt,a4paper]{article}
\usepackage[utf8]{inputenc}
\usepackage[vietnamese]{babel}
\usepackage{amsmath,amssymb,amsfonts}
\usepackage{geometry}
\geometry{a4paper, left=10mm, right=10mm, top=10mm, bottom=10mm}
\usepackage{tikz}
\usetikzlibrary{shapes,arrows.meta,positioning,calc}
\usepackage[most]{tcolorbox}
\usepackage{array,booktabs,multicol}
\usepackage{fontawesome}
\usepackage{helvet}
\renewcommand{\familydefault}{\sfdefault}
\usepackage{parskip}

\definecolor{primary}{RGB}{112, 48, 160}
\definecolor{teal}{RGB}{0, 128, 128}
\definecolor{orange}{RGB}{230, 108, 44}
\definecolor{green}{RGB}{46, 139, 87}
\definecolor{rose}{RGB}{192, 0, 0}

\pagestyle{empty}

\begin{document}

% --- HEADER BANNER ---
\begin{tcolorbox}[
  enhanced,
  colback=violet!8!white,
  colframe=violet!80!black,
  arc=2.5mm, boxrule=1pt,
  left=3mm, right=3mm, top=2mm, bottom=2mm
]
  \begin{tabular*}{\textwidth}{@{\extracolsep{\fill}} lr}
    {\small\bfseries\color{violet!85!black}\faBook\ SỔ TAY KIẾN THỨC TOÁN 2} & 
    {\small\bfseries\color{teal!85!black}\faGraduationCap\ CHƯƠNG TRÌNH GDPT 2018}
  \end{tabular*}
  \vspace{-1mm}
  \begin{center}
    \rule{0.98\textwidth}{0.6pt}
  \end{center}
  \vspace{-2mm}
  \centering
  {\LARGE\bfseries\color{violet!90!black}\faBookmark\ BÀI 66: CHẮC CHẮN, CÓ THỂ, KHÔNG THỂ}\\[1mm]
  {\itshape\footnotesize\color{black!75} Làm quen với các khả năng xảy ra của một sự kiện trong thực tế}
\end{tcolorbox}

\vspace{-2mm}

\noindent
\begin{minipage}[t]{0.49\textwidth}\vspace{0pt}
\begin{tcolorbox}[
  enhanced,
  colback=violet!4!white,
  colframe=violet!75!black,
  arc=2mm, boxrule=0.8pt,
  left=2mm, right=2mm, top=1.8mm, bottom=1.8mm,
  title={\footnotesize\bfseries\color{white}\faKey\ 1. Ý nghĩa của ba cụm từ xác suất}
]
{\footnotesize
\begin{itemize}\setlength{\itemsep}{1pt}\setlength{\leftskip}{-2.5mm}
  \item Khi dự đoán về một sự việc xảy ra trong tương lai hoặc khi lấy ngẫu nhiên:
  \item \textbf{"Chắc chắn":} Sự việc đó \textbf{luôn luôn xảy ra}, $100\%$ không thể khác được.
  \item \textbf{"Có thể":} Sự việc đó \textbf{có khả năng xảy ra hoặc không xảy ra} (chưa biết trước được).
  \item \textbf{"Không thể":} Sự việc đó \textbf{hoàn toàn không bao giờ xảy ra} trong bất kì trường hợp nào.
\end{itemize}
}
\end{tcolorbox}
\end{minipage}
\hfill
\begin{minipage}[t]{0.49\textwidth}\vspace{0pt}
\begin{tcolorbox}[
  enhanced,
  colback=teal!4!white,
  colframe=teal!75!black,
  arc=2mm, boxrule=0.8pt,
  left=2mm, right=2mm, top=1.8mm, bottom=1.8mm,
  title={\footnotesize\bfseries\color{white}\faSliders\ 2. Mô hình thực nghiệm túi bóng màu}
]
{\footnotesize
Trong túi có các quả bóng. Nhắm mắt lấy ra $1$ quả bóng:
\begin{itemize}\setlength{\itemsep}{1pt}\setlength{\leftskip}{-2.5mm}
  \item \textbf{Túi chỉ toàn bóng đỏ:}
  \begin{itemize}\setlength{\itemsep}{0.5pt}\setlength{\leftskip}{-2mm}
    \item Lấy được bóng đỏ: \textbf{Chắc chắn}.
    \item Lấy được bóng xanh: \textbf{Không thể}.
  \end{itemize}
  \item \textbf{Túi có cả bóng đỏ và bóng xanh:}
  \begin{itemize}\setlength{\itemsep}{0.5pt}\setlength{\leftskip}{-2mm}
    \item Lấy được bóng đỏ: \textbf{Có thể}.
    \item Lấy được bóng xanh: \textbf{Có thể}.
    \item Lấy được bóng vàng: \textbf{Không thể} (vì không có bóng vàng).
  \end{itemize}
\end{itemize}
}
\end{tcolorbox}
\end{minipage}

\vspace{-1.5mm}
\begin{tcolorbox}[
  enhanced,
  colback=orange!4!white,
  colframe=orange!80!black,
  arc=2mm, boxrule=0.8pt,
  left=2mm, right=2mm, top=1.5mm, bottom=1.5mm,
  title={\footnotesize\bfseries\color{white}\faSitemap\ 3. Sơ đồ 3 mức độ khả năng xảy ra của một sự kiện}
]
\centering
\begin{tikzpicture}[x=1.65cm, y=0.75cm]
  \tikzset{arr/.style={-{Stealth[scale=0.8]}}, every node/.append style={align=center}}
  \node[rectangle, rounded corners=2pt, draw=green!70!black, fill=green!15, font=\scriptsize\bfseries] (c1) at (0,0) {CHẮC CHẮN\\Xảy ra $100\%$};
  \node[rectangle, rounded corners=2pt, draw=orange!80!black, fill=orange!15, font=\scriptsize\bfseries] (c2) at (2.4,0) {CÓ THỂ\\Có thể xảy ra / không};
  \node[rectangle, rounded corners=2pt, draw=rose!70!black, fill=rose!15, font=\scriptsize\bfseries] (c3) at (4.8,0) {KHÔNG THỂ\\Không bao giờ xảy ra};

  \node[rectangle, rounded corners=2pt, draw=violet!80!black, fill=violet!15, font=\tiny\bfseries] at (2.4,-1) {
    Khả năng phụ thuộc vào điều kiện thực tế của tập hợp đối tượng
  };
  \draw[arr, thick] (c1) |- (2.4,-0.7);
  \draw[arr, thick] (c2) -- (2.4,-0.7);
  \draw[arr, thick] (c3) |- (2.4,-0.7);
\end{tikzpicture}
\end{tcolorbox}

\vspace{-1.5mm}
\noindent
\begin{minipage}[t]{0.49\textwidth}\vspace{0pt}
\begin{tcolorbox}[
  enhanced,
  colback=green!4!white,
  colframe=green!60!black,
  arc=2mm, boxrule=0.8pt,
  left=2mm, right=2mm, top=1.8mm, bottom=1.8mm,
  title={\footnotesize\bfseries\color{white}\faCheckCircle\ 4. Điền cụm từ thích hợp}
]
{\footnotesize
\textbf{1.} Ngày mai mặt trời mọc ở đằng Đông: $\implies$ \textbf{Chắc chắn}.\\
\textbf{2.} Tung đồng xu, mặt sấp xuất hiện: $\implies$ \textbf{Có thể}.\\
\textbf{3.} Một con cá biết bay lượn trên trời như chim: $\implies$ \textbf{Không thể}.\\
\textbf{4.} Trong hộp có $3$ bút chì đỏ và $2$ bút chì xanh. Lấy ra $1$ chiếc bút:
- Lấy được bút đỏ: \textbf{Có thể}.\\
- Lấy được thước kẻ: \textbf{Không thể}.\\
- Lấy được bút viết: \textbf{Chắc chắn}.
}
\end{tcolorbox}
\end{minipage}
\hfill
\begin{minipage}[t]{0.49\textwidth}\vspace{0pt}
\begin{tcolorbox}[
  enhanced,
  colback=rose!4!white,
  colframe=rose!70!black,
  arc=2mm, boxrule=0.8pt,
  left=2mm, right=2mm, top=1.8mm, bottom=1.8mm,
  title={\footnotesize\bfseries\color{white}\faExclamationTriangle\ 5. Mẹo suy luận đúng}
]
{\footnotesize
\begin{itemize}\setlength{\itemsep}{1.5pt}\setlength{\leftskip}{-2.5mm}
  \item Muốn kết luận \textbf{"Chắc chắn"}: Phải đảm bảo tất cả mọi trường hợp đều phải như vậy.
  \item Muốn kết luận \textbf{"Không thể"}: Phải chứng minh hoàn toàn không có đối tượng nào thuộc loại đó.
  \item Nếu có từ $2$ lựa chọn trở lên mà không khẳng định trước được: Ta dùng \textbf{"Có thể"}.
\end{itemize}
}
\end{tcolorbox}
\end{minipage}


\end{document}
